%% toc: P15 160 160.5,160.6,160.7,160.8
% Proofs for the essential-bound consequence (Durrett PTE5 Exercise 3.3.24, p. 143) and Feller's urn identification with explicit weights (Feller II, VII.4, pp. 229--230).
\section{Essential bounds, the moment method and urn parameters}
\label{sec:annex-p15-w02c-p15-012-bounds-and-urns}

\begin{annexitem}{annex:p15-w02c-p15-012-bound-gives-growth}{thm:p15-w02c-p15-012-bound-gives-growth}{Even roots and weak convergence}
\begin{proof}
\emph{Item~1.} Let \(\varepsilon>0\). The number \(\Lambda+\varepsilon\)
exceeds the supremum \(\Lambda\), so it does not belong to
\(\{x:\mathbb P(|X|<x)<1\}\); hence
\(\mathbb P(|X|<\Lambda+\varepsilon)=1\). The events
\(\{|X|<\Lambda+1/j\}\), \(j\geq1\), decrease to \(\{|X|\leq\Lambda\}\),
so continuity from above gives \(\mathbb P(|X|\leq\Lambda)=1\).
Thus \(|X|^k\leq\Lambda^k\) almost surely and
\(\nu_k=\mathbb E|X|^k\leq\Lambda^k<\infty\) for \(k\geq1\).
Since \(|X^k|=|X|^k\), each \(m_k=\mathbb EX^k\) is a finite real number.

\emph{Item~2.} For every \(k\geq1\), \(X^{2k}=|X|^{2k}\), so
\(m_{2k}=\nu_{2k}\geq0\). By \cref{thm:p15-criteria-compact-moments},
\(\nu_j^{1/j}\to\Lambda\) as \(j\to\infty\); the subsequence
\(j=2k\) gives \(m_{2k}^{1/(2k)}\to\Lambda\). By item~1,
\(0\leq m_{2k}^{1/(2k)}=\nu_{2k}^{1/(2k)}\leq\Lambda\), and division by
\(2k\) gives the displayed two-sided bound. The right side tends to
zero, so
\[
 0\leq\liminf_{k\to\infty}\frac{m_{2k}^{1/(2k)}}{2k}
  \leq\limsup_{k\to\infty}\frac{m_{2k}^{1/(2k)}}{2k}
  \leq\lim_{k\to\infty}\frac{\Lambda}{2k}=0 .
\]
Hence the limit superior is \(0\), which is finite.

\emph{Item~3.} The sequence \(m_k\) is a prescribed finite moment sequence
whose normalized even roots have limit superior \(r=0<\infty\).
\Cref{thm:p15-criteria-analytic-determinacy} therefore allows at most one
probability law with these moments. The law of \(X\) has them, so it is
the only one. The hypotheses of \cref{thm:p15-criteria-moment-method}
are: each \(\mu_n\) has all finite absolute moments, its moments converge
to finite real numbers \(m_k\), and the even limits satisfy the growth
condition. The first two are the assumptions placed on \(\mu_n\) in
item~3, and the third is item~2. That theorem gives a unique law \(\mu\)
with moments \(m_k\) and \(\mu_n\Rightarrow\mu\). By the uniqueness just
shown, \(\mu\) is the law of \(X\).

\emph{The uniform example.} For \(X\) uniform on \([0,1]\),
\(\mathbb P(|X|<x)=\min(x,1)\) for \(x\geq0\) and is \(0\) for \(x<0\);
it is below one exactly when \(x<1\), so \(\Lambda=1\). Also
\(m_k=\int_0^1x^k\,dx=1/(k+1)\), and
\(\log m_{2k}^{1/(2k)}=-\log(2k+1)/(2k)\to0\), so
\(m_{2k}^{1/(2k)}\to1\). Fix \(k\geq1\) and write \(f(x)=x^k\), which is
nondecreasing on \([0,1]\). For \(1\leq j\leq n\) and
\(x\in[(j-1)/n,j/n]\), \(f(x)\leq f(j/n)\); integrating over that interval
and summing over \(j\) gives
\(\frac1n\sum_{j=1}^nf(j/n)\geq\int_0^1f=1/(k+1)\). For \(1\leq j\leq n-1\)
and \(x\in[j/n,(j+1)/n]\), \(f(j/n)\leq f(x)\); summing gives
\(\frac1n\sum_{j=1}^{n-1}f(j/n)\leq\int_{1/n}^1f\leq1/(k+1)\), and the
remaining term is \(f(1)/n=1/n\). These are the displayed bounds, so
\(\int x^k\,\mu_n(dx)\to1/(k+1)=m_k\) for every positive integer \(k\).
Each \(\mu_n\) has finite support, so all its absolute moments are
finite, and item~3 shows that \(\mu_n\) converges weakly to the uniform law on \([0,1]\).
\end{proof}
\end{annexitem}

\begin{annexitem}{annex:p15-w02c-p15-012-urn-words}{thm:p15-w02c-p15-012-urn-words}{Products along a word}
\begin{proof}
Fix a word \(x_1,\ldots,x_n\) with \(k\) black entries. Put
\(s_m=x_1+\cdots+x_m\) and
\[
 q_m=\frac{b+cs_m}{b+r+cm}\qquad(0\leq m\leq n-1).
\]
Since \(0\leq s_m\leq m\), \(b\geq0\), \(r\geq0\) and \(c>0\),
\(0\leq b+cs_m\leq b+r+cm\), so \(q_m\in[0,1]\) and
\[
 1-q_m=\frac{r+c(m-s_m)}{b+r+cm}.
\]
Let \(E_m=\{X_1=x_1,\ldots,X_m=x_m\}\), with \(E_0\) the whole space.
By the recursion, \(X_1,\ldots,X_m\) are functions of
\(U_1,\ldots,U_m\), so \(E_m\in\sigma(U_1,\ldots,U_m)\). On \(E_m\),
\(S_m=s_m\) and \(p_m=q_m\). Therefore
\(E_{m+1}=E_m\cap\{U_{m+1}\leq q_m\}\) if \(x_{m+1}=1\) and
\(E_{m+1}=E_m\cap\{U_{m+1}>q_m\}\) if \(x_{m+1}=0\). Since \(U_{m+1}\) is
uniform on \((0,1)\) and independent of \(\sigma(U_1,\ldots,U_m)\),
\[
 \mathbb P(E_{m+1})=\mathbb P(E_m)\,q_m\ \text{ or }\
 \mathbb P(E_m)\,(1-q_m)
\]
in the two cases. No conditional probability given a null event is used.
Multiplying from \(m=0\) to \(n-1\), the denominators are
\(b+r+cm\), \(0\leq m\leq n-1\). At the black steps the numerator is
\(b+cs_m\), and \(s_m\) is the number of earlier black entries, which runs
through \(0,1,\ldots,k-1\) in order. At the red steps the numerator is
\(r+c(m-s_m)\), and \(m-s_m\) is the number of earlier red entries, which
runs through \(0,1,\ldots,n-k-1\). This gives the first displayed
formula, regardless of how the black and red steps are interleaved.
Writing each factor \(a+ic\) as \(c(a/c+i)\), the numerator and
denominator each contain \(n\) factors of \(c\), which cancel; this gives
the rising-product form. The case \(n=0\) is the empty word, of
probability one.

For a permutation \(\pi\) of \(\{1,\ldots,n\}\) and a word \(x\),
\[
 \mathbb P\big((X_{\pi1},\ldots,X_{\pi n})=x\big)
   =\mathbb P\big((X_1,\ldots,X_n)=y\big),\qquad y_{\pi j}=x_j ,
\]
and \(y\) has the same number of black entries as \(x\). By the formula
both sides equal \(\mathbb P((X_1,\ldots,X_n)=x)\). So every finite
vector has a permutation-invariant law. For a permutation of
\(\mathbb N\) moving finitely many coordinates, choose \(N\) at least as
large as every moved coordinate and every coordinate of a given cylinder
event; summing the length-\(N\) word probabilities over the unspecified
coordinates shows that the cylinder has the same probability before and
after the permutation. Cylinder events form a pi-system generating the
product sigma-field of \(\{0,1\}^{\mathbb N}\), so the two laws agree on
that sigma-field. Thus the sequence is an infinite exchangeable binary
sequence. Taking the all-black word, \(k=n\), gives the formula for
\(c_n\).
\end{proof}
\end{annexitem}

\begin{annexitem}{annex:p15-w02c-p15-012-urn-beta}{thm:p15-w02c-p15-012-urn-beta}{The beta integral recursion and Hausdorff uniqueness}
\begin{proof}
\emph{Finiteness.} Let \(\alpha,\beta>0\). On \((0,1/2]\),
\((1-\theta)^{\beta-1}\leq\max(1,2^{1-\beta})\), and
\(\int_0^{1/2}\theta^{\alpha-1}\,d\theta=2^{-\alpha}/\alpha<\infty\)
because \(\alpha>0\). On \([1/2,1)\),
\(\theta^{\alpha-1}\leq\max(1,2^{1-\alpha})\), and
\(\int_{1/2}^1(1-\theta)^{\beta-1}\,d\theta=2^{-\beta}/\beta<\infty\)
because \(\beta>0\). So \(B(\alpha,\beta)<\infty\); it is positive because
its integrand is positive on \((0,1)\). This is where both \(b>0\) and
\(r>0\) are used: \(\alpha=b/c\) and \(\beta=r/c\) must both be positive.

\emph{Recursion.} For \(0<\varepsilon<1/2\), integration by parts with
\(u=(1-\theta)^\beta\) and \(v=\theta^\alpha/\alpha\) on
\([\varepsilon,1-\varepsilon]\) gives
\[
 \int_\varepsilon^{1-\varepsilon}\theta^{\alpha-1}(1-\theta)^\beta\,d\theta
 =\Big[\frac{\theta^\alpha(1-\theta)^\beta}{\alpha}\Big]_\varepsilon^{1-\varepsilon}
  +\frac\beta\alpha\int_\varepsilon^{1-\varepsilon}
        \theta^\alpha(1-\theta)^{\beta-1}\,d\theta .
\]
As \(\varepsilon\downarrow0\), the bracket tends to zero because
\(\alpha,\beta>0\), and both integrals of nonnegative functions increase
to \(B(\alpha,\beta+1)\) and \(B(\alpha+1,\beta)\), which are finite by
the first paragraph. Hence \(B(\alpha,\beta+1)=(\beta/\alpha)B(\alpha+1,\beta)\).
Since \(\theta+(1-\theta)=1\),
\(\theta^{\alpha-1}(1-\theta)^{\beta-1}
 =\theta^\alpha(1-\theta)^{\beta-1}+\theta^{\alpha-1}(1-\theta)^\beta\)
on \((0,1)\), so
\(B(\alpha,\beta)=B(\alpha+1,\beta)+B(\alpha,\beta+1)
 =B(\alpha+1,\beta)(\alpha+\beta)/\alpha\). Therefore
\[
 B(\alpha+1,\beta)=\frac{\alpha}{\alpha+\beta}B(\alpha,\beta),\qquad
 B(\alpha,\beta+1)=\frac{\beta}{\alpha+\beta}B(\alpha,\beta).
\]
Applying the first identity \(n\) times, with first arguments
\(\alpha,\alpha+1,\ldots,\alpha+n-1\), gives
\(B(\alpha+n,\beta)/B(\alpha,\beta)=(\alpha)_n/(\alpha+\beta)_n\).

\emph{Item~1.} The \(n\)-th moment of \(\operatorname{beta}(\alpha,\beta)\)
is \(B(\alpha,\beta)^{-1}\int_0^1\theta^{\alpha+n-1}(1-\theta)^{\beta-1}d\theta
=B(\alpha+n,\beta)/B(\alpha,\beta)=(\alpha)_n/(\alpha+\beta)_n\).
For \(y>0\), integration by parts on \([\varepsilon,R]\) gives
\(\int_\varepsilon^Rt^ye^{-t}dt=\varepsilon^ye^{-\varepsilon}-R^ye^{-R}
+y\int_\varepsilon^Rt^{y-1}e^{-t}dt\); letting \(\varepsilon\downarrow0\)
and \(R\to\infty\) gives \(\Gamma(y+1)=y\Gamma(y)\). Iterating,
\(\Gamma(x+n)=(x)_n\Gamma(x)\) for \(x>0\), so
\((\alpha)_n/(\alpha+\beta)_n
 =\Gamma(\alpha+n)\Gamma(\alpha+\beta)/(\Gamma(\alpha+\beta+n)\Gamma(\alpha))\).
By \cref{thm:p15-w02c-p15-012-urn-words}, with \(b/c=\alpha\) and
\((b+r)/c=\alpha+\beta\), this is \(c_n\). The case \(n=0\) gives
\(c_0=1\), the total mass.

\emph{Item~2.} By \cref{thm:p15-w02c-p15-012-urn-words} the draws form an
infinite exchangeable binary sequence, so
\cref{thm:p15-feller-binary-mixture} applies. It gives a probability
\(\lambda\) on \([0,1]\) with
\(p_{k,n}=\int\theta^k(1-\theta)^{n-k}\,\lambda(d\theta)\); the case
\(k=n\) reads \(c_n=\int\theta^n\,\lambda(d\theta)\). By item~1,
\(\operatorname{beta}(\alpha,\beta)\) is also a probability on \([0,1]\)
with moments \(c_n\) for every \(n\geq0\). By \cref{thm:p15-feller-hausdorff},
a sequence is the moment sequence of at most one probability on \([0,1]\);
hence \(\lambda=\operatorname{beta}(\alpha,\beta)\).

\emph{Item~3.} Applying the first recursion \(k\) times with second
argument \(\beta+n-k\), then the second recursion \(n-k\) times,
\begin{align*}
 B(\alpha+k,\beta+n-k)
   &=\frac{(\alpha)_k}{(\alpha+\beta+n-k)_k}B(\alpha,\beta+n-k)\\
   &=\frac{(\alpha)_k\,(\beta)_{n-k}}
         {(\alpha+\beta+n-k)_k\,(\alpha+\beta)_{n-k}}B(\alpha,\beta).
\end{align*}
The two rising products in the denominator multiply to
\((\alpha+\beta)_n\). So
\(B(\alpha+k,\beta+n-k)/B(\alpha,\beta)=(\alpha)_k(\beta)_{n-k}/(\alpha+\beta)_n\),
which is the word probability of \cref{thm:p15-w02c-p15-012-urn-words}
with \(r/c=\beta\). It is also
\(\int\theta^k(1-\theta)^{n-k}\,\operatorname{beta}(\alpha,\beta)(d\theta)\),
in agreement with item~2. The event \(\{S_n=k\}\) is the disjoint union
of the \(\binom nk\) word events with \(k\) black entries, each of this
probability, which gives the count law.
\end{proof}
\end{annexitem}

\begin{annexitem}{annex:p15-w02c-p15-012-urn-one-colour}{thm:p15-w02c-p15-012-urn-one-colour}{The one-colour urns}
\begin{proof}
\emph{Item~1, \(b>0=r\).} We show by induction on \(m\) that
\(X_1=\cdots=X_m=1\) at every sample point. For \(m=0\) there is nothing
to prove. If it holds for \(m\), then \(S_m=m\) and
\(p_m=(b+cm)/(b+cm)=1\); since \(U_{m+1}<1\), \(X_{m+1}=1\).
Hence \(c_n=1\) for every \(n\), and the constant sequence is
exchangeable. The product formula of
\cref{thm:p15-w02c-p15-012-urn-words} with \(k=n\) and \(r=0\) has equal
numerator and denominator, and \((b/c)_n/(b/c)_n=1\).
\Cref{thm:p15-feller-binary-mixture} gives a probability \(\lambda\) on
\([0,1]\) with moments \(c_n=1\). The point mass \(\delta_1\) has moments
\(1^n=1\), so \cref{thm:p15-feller-hausdorff} gives \(\lambda=\delta_1\).
With \(\alpha=b/c>0\), on \([1/2,1)\) we have
\(\theta^{\alpha-1}\geq\min(1,2^{1-\alpha})>0\), so
\(\int_{1/2}^1\theta^{\alpha-1}(1-\theta)^{-1}d\theta
 \geq\min(1,2^{1-\alpha})\int_{1/2}^1(1-\theta)^{-1}d\theta=\infty\);
\(B(\alpha,0)\) diverges and no normalized beta density exists.

\emph{Item~2, \(b=0<r\).} By induction, \(X_1=\cdots=X_m=0\) at every
sample point: if it holds for \(m\), then \(S_m=0\),
\(p_m=0/(r+cm)=0\), and \(U_{m+1}>0\) gives \(X_{m+1}=0\). Hence
\(c_n=\mathbb P(X_1=1,\ldots)=0\) for \(n\geq1\), and the first factor
\(b=0\) of the product formula gives the same value. The point mass
\(\delta_0\) has moments \(0^n=0\) for \(n\geq1\) and \(1\) for \(n=0\);
the binary mixture theorem and Hausdorff uniqueness give
\(\lambda=\delta_0\). With \(\beta=r/c>0\), on \((0,1/2]\),
\((1-\theta)^{\beta-1}\geq\min(1,2^{1-\beta})>0\), so
\(\int_0^{1/2}\theta^{-1}(1-\theta)^{\beta-1}d\theta=\infty\) and
\(B(0,\beta)\) diverges. Finally
\(\Gamma(x)=\Gamma(x+1)/x\) for \(x>0\), and \(\Gamma(x+1)\to\Gamma(1)=1\)
as \(x\downarrow0\) by dominated convergence with the integrable bound
\(t^xe^{-t}\leq(1+t)e^{-t}\) for \(0<x\leq1\); hence
\(\Gamma(x)\to\infty\) as \(x\downarrow0\), and the gamma expression,
which contains \(\Gamma(b/c)=\Gamma(0)\), is undefined.
\end{proof}
\end{annexitem}
